\documentclass[conference]{IEEEtran}
\usepackage[utf8]{inputenc}
\usepackage[T1]{fontenc}
\usepackage{graphicx}
\usepackage{amsmath}
\usepackage{booktabs}
\usepackage{hyperref}
\usepackage[section]{placeins}

\title{Unsupervised Neural Network for Multi-Genre Music Generation}
\author{\IEEEauthorblockN{Abrar Fahim}
\IEEEauthorblockA{\textit{Department of Computer Science} \\
\textit{BRAC University}\\
Student ID: 22201689}}

\begin{document}
\maketitle

\begin{abstract}
This report presents a four-stage development of multi-genre (Classical, Jazz, Rock, Pop and Electronic) music generation pipeline using unsupervised and preference-aligned learning: (1) LSTM autoencoder, (2) multi-genre VAE, (3) autoregressive Transformer and (4) RLHF fine-tuning. Their performances are evaluated thoroughly with quantitative metrics such as histogram similarity, rhythm diversity scores and repetition ratios. The generative capabilities of the above mentioned models are compared against naive random generators and Markov Chain models to contrast their capabilities.
\end{abstract}

\section{Introduction}
Music creation is a significant difficulty in artificial intelligence because it involves both localized structural coherence and long-term musical connection. Symbolic music, particularly in the form of MIDI data, provides a discrete and highly organized representation of musical works that are free of audio waveform complexity. The primary goal of this project is to develop, train and formally evaluate more complex unsupervised generative models in symbolic music sequences. 

The analysis includes many architectural milestones trained on datasets like MAESTRO and Lakh MIDI. Initially, we create basic baselines using Random Note Selection and Markov Chains. Following that, we use an LSTM Autoencoder developed specifically for sequence reconstruction and representation learning (Task 1). Building on this framework, we develop a Variational Autoencoder (VAE) that enforces a regularized latent space, facilitating the variety and latent interpolation required for multi-genre symbolic generation (Task 2). To address the limits of recurrent architectures in terms of durable long-term structures, we then propose a decoder-only Transformer model (Task 3). Finally, we investigate how to align our autoregressive models with nuanced human preferences using Reinforcement Learning from Human Feedback (RLHF) techniques (Task 4). This study describes the methods, experimental settings and performance evaluations for various generative paradigms.

\section{Related Work}
Algorithmic music creation has evolved through multiple generations of technical innovation. Early stochastic models, such as Markov chains, gained popularity because of their ability to accurately describe changes between notes. However, these localized statistical models were unable to preserve overall consistent melodic phrasing throughout the long runs.

Later advancements in deep recurrent neural networks, particularly Long Short-Term Memory (LSTM) networks, enabled models to avoid the vanishing gradient problem and remember longer temporal contexts. Autoencoders using LSTM layers have been proven excellent at compressing symbolic music into fixed-dimensional latent vectors and then recreating them. The Variational Autoencoder (VAE) introduces generative variety and enables continuous interpolation between discrete musical motifs by regularizing the learned distribution to a Gaussian prior.

More recently, attention-based architectures, notably the Transformer, have significantly improved their ability to imitate the musical composition with their long-span dependencies. These Transformer models rely primarily on self-attention processes to create music with continuous themes and rhythmic patterns that have extended duration compared to previous technologies. Further advancements include fine-tuning a pre-trained generative model with Reinforcement Learning using Human Feedback (RLHF). This is to optimize output generation against subjective perceptual qualities, a domain well-studied in large language models but recently extended to abstract symbolic modalities. The idea is to let actual humans judge what is a good output and what is a bad output. Then based on their reviews we model a reward function that let's the generative model know which patterns are discouraged to follow.
\section{Methodology}

\subsection{Task 1: LSTM Autoencoder}
The first generative model instituted in this study is a recurrent sequence-to-sequence Autoencoder. The encoder is built using an LSTM network that processes tokenized sequential inputs derived from piano roll matrices and MIDI event tokens into a compressed fixed-length hidden bottleneck vector $z$. The decoder, symmetrically structured as an LSTM, is initialized with $z$ as its initial hidden state. It attempts to reconstruct the input sequence autoregressively.

The network is trained with the aim to minimize reconstruction loss (e.g., Mean Squared Error for continuous representations or Cross-Entropy for categorical tokens). Trying to compel the encoder to isolate the most salient structural attributes of the music. However, a vanilla autoencoder fundamentally lacks the probabilistic constraints when it tries to represent a latent space. Standardizing this generation process necessitates fitting a multi-variate Gaussian distribution which is defined by the empirical mean and standard deviation over the entire training corpus before sampling random $z$ vectors. 

\subsection{Task 2: Variational Autoencoder (VAE)}
To address the latent discontinuities that characterize the standard autoencoder, we then move on to implement a Variational Autoencoder (VAE). This change causes a modification of the output parameters by dictating a probability distribution. Specifically, a mean $\mu$ and logarithmic variance $\log(\sigma^2)$ for each latent dimension. Using this reparameterization trick, the decoder then executes its reconstructions upon sampling from this learned distribution.

The objective function simultaneously optimizes for both reconstruction fidelity and Kullback-Leibler (KL) divergence. This regularizes the encoded distributions towards a standard normal prior $\mathcal{N}(0, I)$. This regularization enforces a smooth and densely packed latent topology, which allows us to produce musical morphs between disparate genres that are cohesive by performing linear interpolations between discrete latent vectors.

\subsection{Task 3: Autoregressive Transformer}
Recurrent networks which process tokens sequentially face a limitation in capturing the global dependencies in extended musical pieces. Hence, we design a decoder-only autoregressive Transformer model. The inputs are discretely embedded and injected with positional encodings to ensure temporal awareness. 

The core of the architecture deploys stacked layers of masked multi-head self-attention. By restricting attention to strictly preceding tokens (causal masking), the network learns to accurately estimate the conditional probability of subsequent musical events given the unrolled context. Trained using a standard negative log-likelihood paradigm (Cross-Entropy Loss), performance is principally assessed via test-set perplexity. For generative application, the model recursively samples from its output softmax distribution utilizing top-k/top-p filtering or temperature scaling to balance stochastic creativity versus structural rigidity over prolonged sequences.

\subsection{Task 4: RLHF Fine-Tuning}
For Task 4, we fine-tune the Transformer policy using reinforcement learning from human feedback (RLHF). Human survey ratings are used to train a reward model over symbolic-music quality features and policy updates are then performed to maximize expected reward under this learned preference signal. The training process logs mean reward and policy loss over iterative optimization steps and final evaluation compares pre-RLHF and post-RLHF generations with objective metrics; survey-based post-RLHF comparison is included when available in exported metadata.

\subsection{Architecture Details and Hyperparameters}
To keep experiments reproducible, Table~\ref{tab:arch-hparams} summarizes the exact architecture and optimization settings used in code. Recurrent models (Task 1 and Task 2) operate on piano-roll tensors of shape $(64,128)$, while Transformer/RLHF stages (Task 3 and Task 4) operate on event-token sequences with maximum length 256.

\begin{table}[!t]
\caption{Model Architecture and Training Hyperparameters}
\label{tab:arch-hparams}
\begin{center}
\begin{tabular}{lcc}
\toprule
Setting & Value & Applied In \\
\midrule
LSTM hidden size & 256 & Task 1, Task 2 \\
Latent dimension & 128 & Task 1, Task 2 \\
AE batch size / epochs / lr & 64 / 10 / $10^{-3}$ & Task 1 \\
VAE $\beta$ (KL weight) & 1.0 & Task 2 \\
VAE batch size / epochs / lr & 64 / 10 / $10^{-3}$ & Task 2 \\
Transformer $d_{model}$ & 256 & Task 3, Task 4 policy \\
Transformer heads / layers & 8 / 6 & Task 3, Task 4 policy \\
Transformer max sequence length & 256 tokens & Task 3, Task 4 policy \\
Transformer batch size / lr & 64 / $10^{-4}$ & Task 3 \\
RLHF steps / episodes per step & 200 / 16 & Task 4 run \\
RLHF policy lr & $5\times10^{-6}$ & Task 4 \\
RLHF top-$k$ / temperature & 32 / 1.0 & Task 4 \\
Reward-model feature dimension & 8 & Task 4 \\
\bottomrule
\end{tabular}
\end{center}
\end{table}

Task 1 and Task 2 share an encoder-decoder LSTM backbone, but Task 2 replaces deterministic encoding with probabilistic encoding $(\mu,\log\sigma^2)$ and reparameterized latent sampling. Task 3 introduces masked self-attention with causal decoding to model long-range event dependencies. Task 4 reuses the Task 3 Transformer as the policy network and updates it through REINFORCE using survey-conditioned reward signals.

\begin{figure}[!t]
\centering
\small
\begin{tabular}{|c|c|}
\hline
\textbf{Task 1 (AE)} & \textbf{Task 2 (VAE)} \\
\hline
\begin{minipage}[t]{0.45\columnwidth}\raggedright
Input piano-roll $(64,128)$ \newline $\rightarrow$ LSTM encoder (256) \newline $\rightarrow$ latent $z$ (128) \newline $\rightarrow$ LSTM decoder (256) \newline $\rightarrow$ reconstructed output
\end{minipage} &
\begin{minipage}[t]{0.45\columnwidth}\raggedright
Input piano-roll $(64,128)$ \newline $\rightarrow$ LSTM encoder \newline $\rightarrow$ $(\mu,\log\sigma^2)$ (128) \newline $\rightarrow$ reparameterization \newline $\rightarrow$ LSTM decoder $\rightarrow$ output
\end{minipage} \\
\hline
\textbf{Task 3 (Transformer)} & \textbf{Task 4 (RLHF)} \\
\hline
\begin{minipage}[t]{0.45\columnwidth}\raggedright
Event tokens (max 256) + positional encoding \newline $\rightarrow$ 6 masked self-attention \newline layers ($d_{model}=256$, 8 heads) \newline $\rightarrow$ next-token logits
\end{minipage} &
\begin{minipage}[t]{0.45\columnwidth}\raggedright
Task 3 policy \newline $\rightarrow$ rollout generation \newline $\rightarrow$ reward model \newline $\rightarrow$ REINFORCE update \newline $\rightarrow$ tuned policy
\end{minipage} \\
\hline
\end{tabular}
\caption{Compact architecture schematics for Tasks 1--4.}
\label{fig:task-architectures}
\normalsize
\end{figure}

\subsection{Computational Setup and Efficiency}
All experiments were executed on a desktop workstation with an Intel ARC B580 GPU, AMD RYZEN 7500F CPU and 16GB DDR5 RAM. This hardware profile is important for interpreting both throughput and feasible model scale under the project timeline.

For data preparation, the largest memory pressure came from building and storing the Lakh-derived processed splits (train/validation/test) as dense numpy arrays before minibatch loading. The 64GB system RAM was sufficient to complete preprocessing and split persistence without external sharding, while keeping iterative experimentation practical.

For model training, the main bottlenecks appeared in Task 3 and Task 4: (i) masked self-attention over longer token contexts in the Transformer and (ii) repeated sequence rollout plus log-probability recomputation in RLHF. The implementation mitigates memory overhead by casting training batches to \texttt{float32} on-device, limiting sequence length to 256 tokens and using staged RLHF updates (rollout, advantage normalization, then backpropagation) so that gradients are retained only when needed.

Although wall-clock epoch timing was not explicitly profiled in this run, observed behavior was consistent with expected scaling: preprocessing remained CPU/RAM bound, while Transformer and RLHF optimization were GPU-memory sensitive and benefited from conservative batch and sequence configurations.

\section{Dataset and Preprocessing}
The project uses two symbolic MIDI corpora with distinct roles. MAESTRO is used for Task 1 (single-domain reconstruction), while Lakh is used for Task 2--Task 4 (multi-genre generation and alignment). The processed split sizes in this run are:

\begin{itemize}
\item MAESTRO train/validation/test: 143,037 / 17,455 / 17,948 segments.
\item Lakh train/validation/test: 799,444 / 99,930 / 99,932 segments.
\end{itemize}

Each segment is stored as a fixed-size piano-roll tensor $(64,128)$ in \texttt{uint8}. Here, 64 denotes temporal steps and 128 denotes MIDI pitches. MIDI is discretized at sampling frequency $\mathrm{FS}=16$ and binarized per timestep-pitch cell. During training, batches are cast to \texttt{float32} on-device to reduce host-memory overhead.

For Task 3 and Task 4, piano-rolls are converted into event-token sequences (max length 256). Token categories are PAD/BOS/EOS, note-on (128 IDs), note-off (128 IDs), velocity bins (32 IDs), duration tokens (64 IDs) and time-shift tokens (32 IDs), giving total vocabulary size $3+128+128+32+64+32=387$. Long durations and long silences are represented via chunked duration/time-shift events up to configured maxima.

\begin{figure}[htbp]
\centering
\footnotesize
\begin{tabular}{|l|}
\hline
\textbf{End-to-End Project Pipeline} \\
\hline
MIDI datasets (MAESTRO/Lakh) $\rightarrow$ preprocessing and split generation \\
$\rightarrow$ Task 1 (LSTM-AE) $\rightarrow$ Task 2 (VAE) \\
$\rightarrow$ Task 3 (Transformer) $\rightarrow$ Task 4 (RLHF over Transformer) \\
$\rightarrow$ objective metrics and human evaluation tables \\
\hline
\end{tabular}
\caption{End-to-end project pipeline architecture.}
\label{fig:pipeline-overview}
\normalsize
\end{figure}

\subsection{Tokenization Strategy}
To process the sparse $(64, 128)$ piano-roll matrices, we implement a custom event-based tokenizer with a compact vocabulary of 387 tokens. This vocabulary comprises control tokens (PAD, BOS, EOS), note events (128 Note-On, 128 Note-Off), 32 quantized velocity bins and temporal tokens (64 Duration, 32 Time-Shift).

Sequences are structured using a strict grammatical pattern: \texttt{[time\_shift] $\rightarrow$ note\_on $\rightarrow$ velocity $\rightarrow$ duration $\rightarrow$ [note\_off]}. Extended notes or silences exceeding maximum token bounds are algorithmically decomposed into maximum-sized sequential chunks. To optimize memory usage during Transformer training, all sequences are strictly truncated or padded to a length of 256 tokens.

During inference, the decoder utilizes heuristic parsing to discard hallucinated or orphaned tokens (e.g., a duration without a preceding note event), ensuring fault-tolerant reconstruction back to a $(128, T)$ piano-roll tensor. This bounded vocabulary and fixed sequence length constrain the quadratic memory complexity of masked self-attention. Consequently, this design ensures that training batches of size 64 fit comfortably within the VRAM of an NVIDIA RTX 4080 Super, while reserving the 64GB DDR5 system memory for caching the large-scale Lakh MIDI dataset splits.

\section{Evaluation Metrics}
To evaluate generated symbolic music, we use three objective metrics and one subjective metric:

\textbf{Pitch Histogram Similarity}: histogram distance between generated and reference pitch distributions,
\begin{equation}
H(p,q)=\sum_{i=1}^{12}|p_i-q_i|.
\end{equation}
and similarity is computed as
\begin{equation}
S_{\text{pitch}} = 1 - 0.5\,H(p,q),\quad S_{\text{pitch}}\in[0,1],
\end{equation}
where higher is better.

\textbf{Rhythm Diversity Score}: ratio of unique note durations to total notes,
\begin{equation}
D_{\text{rhythm}}=\frac{\#\text{unique durations}}{\#\text{total notes}}.
\end{equation}

\textbf{Repetition Ratio}: fraction of repeated local patterns,
\begin{equation}
R=\frac{\#\text{repeated patterns}}{\#\text{total patterns}}.
\end{equation}

\textbf{Human Listening Score}: mean survey rating on a $1$--$5$ scale,
\begin{equation}
\text{Score}_{\text{human}} \in [1,5].
\end{equation}

\section{Experiments and Results}
This section reports completed training and evaluation evidence for all four tasks. Quantitative metrics are reported from exported CSV files and summarized in IEEE-style tables.

\subsection{Deliverable}
All task deliverables required in the assignment are completed:
\begin{itemize}
\item Task 1: autoencoder implementation, reconstruction loss curve and 5 generated MIDI samples.
\item Task 2: VAE with KL-divergence objective, 8 generated samples and 8 latent interpolation samples.
\item Task 3: Transformer implementation, perplexity report and 10 generated long-sequence compositions.
\item Task 4: listening survey data (10 participants per surveyed sample), reward-model based RLHF run, 10 RLHF output samples and before/after comparison.
\end{itemize}

\subsection{Baseline and Task-Level Comparison}
Table~\ref{tab:baseline-task-comparison} replicates the assignment-style performance comparison format across baseline and task models.

\begin{table*}[!t]
\caption{Performance Comparison Table (Task 1--Task 4)}
\label{tab:baseline-task-comparison}
\begin{center}
\begin{tabular}{lccccc}
\toprule
Model & Loss & Perplexity & Rhythm Diversity & Human Score & Genre Control \\
\midrule
Random Generator & -- & -- & 0.0235 & 3.10 & None \\
Markov Chain & -- & -- & 0.5912 & 3.75 & Weak \\
Task 1: Autoencoder & -- & -- & 0.8862 & 3.85 & Single Genre \\
Task 2: VAE Multi-Genre & -- & -- & 0.0860 & 3.18 & Moderate \\
Task 3: Transformer & -- & 2.6312 (best val) & 0.3452 & 3.13 & Strong \\
Task 4: RLHF-Tuned Model & -- & -- & 0.5777 & -- & Strong \\
\bottomrule
\end{tabular}
\end{center}
\end{table*}

\subsection{Objective Evaluation Across Models}
Table~\ref{tab:all-models-objective} summarizes the objective metrics from the consolidated evaluation output.

\begin{table}[!t]
\caption{Objective Metric Comparison Across Models}
\label{tab:all-models-objective}
\begin{center}
\begin{tabular}{lccc}
\toprule
Model & Rhythm Div. & Repetition & Pitch Hist. Sim. \\
\midrule
Markov baseline & 0.5912 & 0.0000 & 0.6268 \\
Random baseline & 0.0235 & 0.0000 & 0.8896 \\
Task 1 (AE) & 0.8862 & 0.0000 & 0.4476 \\
Task 2 (VAE) & 0.0860 & 0.0000 & 0.7965 \\
Task 3 (Transformer) & 0.3452 & 0.0577 & 0.5364 \\
Task 4 (RLHF after) & 0.5777 & 0.0000 & 0.5298 \\
\bottomrule
\end{tabular}
\end{center}
\end{table}

\subsection{Qualitative Results: Piano-roll Structure}
To complement objective metrics, Figure~\ref{fig:qualitative-pianorolls} compares representative piano-roll snippets from three model classes: Random baseline, Task 1 Autoencoder and Task 4 RLHF outputs.

\begin{figure}[htbp]
\centerline{\includegraphics[width=0.98\columnwidth]{./outputs/plots/qualitative_pianoroll_comparison.png}}
\caption{Qualitative piano-roll comparison across Random baseline, Task 1 (AE) and Task 4 (RLHF).}
\label{fig:qualitative-pianorolls}
\end{figure}

The Random baseline exhibits scattered activations with weak local continuity, which visually corresponds to low compositional structure. Task 1 shows stronger note occupancy and short-range coherence, but also repeated texture bands that suggest limited long-horizon variation. Task 4 RLHF displays clearer phrase-like groupings and more controlled rhythmic density across adjacent windows, consistent with the reduction in repetition ratio and improved human-score trend reported in Tables~\ref{tab:task4-objective-compare} and \ref{tab:task4-survey-compare}. While these snapshots are not a substitute for listening tests, they provide a useful structural sanity check alongside numerical metrics.

\subsection{Task 1: Autoencoder Training Dynamics}
Figure~\ref{fig:task1-loss} shows the Task 1 training curve. The reconstruction loss demonstrates stable optimization behavior and supports that the latent bottleneck learns compressive representations that are sufficient for sequence reconstruction.

\begin{figure}[htbp]
\centerline{\includegraphics[width=0.95\columnwidth]{./outputs/plots/task1_loss.png}}
\caption{Task 1 LSTM autoencoder loss curve.}
\label{fig:task1-loss}
\end{figure}

\subsection{Task 2: VAE Training Dynamics}
Figure~\ref{fig:task2-loss} presents the Task 2 VAE loss trend. The combined objective (reconstruction + KL regularization) indicates that the model is learning a structured latent space while maintaining acceptable reconstruction quality.

\begin{figure}[htbp]
\centerline{\includegraphics[width=0.95\columnwidth]{./outputs/plots/task2_vae_loss.png}}
\caption{Task 2 VAE training curve.}
\label{fig:task2-loss}
\end{figure}

In addition to random latent sampling, Task 2 includes a latent interpolation experiment using 8 intermediate points (\texttt{task2\_interp\_0.mid} to \texttt{task2\_interp\_7.mid}). The interpolation set shows stable objective behavior (mean rhythm diversity $0.1053$, repetition $0.0000$, pitch-histogram similarity $0.7862$), indicating continuity in decoded musical structure across neighboring latent positions. Compared with random Task 2 samples (mean rhythm diversity $0.0860$, pitch-histogram similarity $0.7965$), interpolation preserves tonal proximity while producing smoother transitions, which is the expected behavior of a regularized latent space.

\subsection{Task 3: Transformer Training}
Figure~\ref{fig:task3-curves} shows the full Task 3 training dynamics across 50 cumulative epochs. The model converges to a best validation loss of 0.9674 and a best validation perplexity of 2.6312, with last-epoch validation perplexity 2.6349, indicating stable autoregressive next-token modeling after extended training.

\begin{table}[!t]
\caption{Task 3 Transformer Summary}
\label{tab:task3-summary}
\begin{center}
\begin{tabular}{lc}
\toprule
Metric & Value \\
\midrule
Total epochs trained & 50 \\
Best validation loss & 0.9674 \\
Best validation perplexity & 2.6312 \\
Last validation loss & 0.9688 \\
Last validation perplexity & 2.6349 \\
\bottomrule
\end{tabular}
\end{center}
\end{table}

\begin{figure}[htbp]
\centerline{\includegraphics[width=0.95\columnwidth]{./outputs/plots/task3_transformer_curves.png}}
\caption{Task 3 Transformer training and validation curves after full training.}
\label{fig:task3-curves}
\end{figure}

\subsection{Task 4: RLHF Training and Evaluation}
RLHF optimization completed with survey-model reward as the training signal. The cumulative training log spans 200 optimization steps trained over time with 50 steps each and  16 episodes per step. The best observed mean reward was 0.7232 and the final mean reward was 0.7232. The companion summary JSON records \texttt{rl\_steps=50}, which corresponds to a checkpoint summary rather than the full cumulative log length.

Survey protocol: each surveyed sample was scored by 10 participants on a 1--5 scale. A total of 41 files are recorded in the survey metadata used for reward-model fitting and summary analysis.

The RLHF optimization loop is implemented in three phases per step: (1) rollout generation without gradient tracking to collect sampled sequences and rewards, (2) reward standardization and advantage computation and (3) per-episode log-probability recomputation with gradients followed by policy updates. This staged design improves memory behavior while preserving policy-gradient correctness on long symbolic sequences.

Reward learning includes a dedicated 8-feature MLP reward model trained from survey labels. Features are: pitch similarity, rhythm diversity, anti-repetition, note density, duration variance, mean velocity, pitch span and inter-onset regularity. The model architecture is \texttt{Linear(8,64)-ReLU-\linebreak Linear(64,32)-ReLU-Linear(32,1)-Sigmoid}. During RLHF, this learned reward is preferred; fallback paths combine proxy reward and survey mean when direct lookup is unavailable.

\paragraph{Feature Justification.}
The reward features were selected to balance tonal, rhythmic and stylistic criteria relevant to multi-genre symbolic music. Pitch similarity anchors generated material to plausible tonal distributions; rhythm diversity discourages metrical collapse into a single duration pattern; and anti-repetition penalizes short-loop artifacts common in autoregressive decoding. Note density and duration variance jointly regulate texture thickness and temporal expressiveness, preventing outputs that are either empty or uniformly saturated. Mean velocity contributes dynamic control, pitch span captures melodic range and inter-onset regularity provides a compact descriptor of rhythmic smoothness versus jitter.

\paragraph{Synthetic Preference Protocol.}
Fully manual human-in-the-loop labeling at RLHF scale was outside the scope of this course timeline, so survey labels were synthesized from proxy quality signals with participant-specific perturbations. This design allows controlled multi-rater behavior while preserving a realistic spread of opinions through additive Gaussian noise and per-rater bias terms. As a result, the reward model is trained on noisy but structured preference targets rather than deterministic scores, which better approximates inter-rater disagreement than a single-point labeling scheme.

\begin{figure}[htbp]
\centerline{\includegraphics[width=0.95\columnwidth]{./outputs/plots/task4_rlhf_curves.png}}
\caption{Task 4 RLHF reward and policy-loss curves.}
\label{fig:task4-rlhf-curves}
\end{figure}

\begin{table}[!t]
\caption{Task 4 Objective Comparison (Before vs. After RLHF)}
\label{tab:task4-objective-compare}
\begin{center}
\begin{tabular}{lccc}
\toprule
Group & Rhythm Div. & Repetition & Pitch Hist. Sim. \\
\midrule
Before (Task 3) & 0.3452 & 0.0577 & 0.5364 \\
After (Task 4) & 0.5777 & 0.0000 & 0.5298 \\
\bottomrule
\end{tabular}
\end{center}
\end{table}

\begin{table*}[!t]
\caption{Task 4 Survey Comparison (Before vs. After RLHF)}
\label{tab:task4-survey-compare}
\begin{center}
\begin{tabular}{lccc}
\toprule
Group & Mean Human Score & Mean Proxy Reward & Avg. Human Std \\
\midrule
Before (Task 3) & 3.1300 & 0.5411 & 0.4335 \\
After (Task 4) & 3.5300 & 0.6430 & 0.4120 \\
\bottomrule
\end{tabular}
\end{center}
\end{table*}

The results indicate favorable changes in the RLHF outputs relative to the Task 3 baseline on this run: higher rhythm diversity and reduced repetition, with a slight decrease in pitch-histogram similarity (0.5364 $\rightarrow$ 0.5298). The regenerated survey metadata now also shows improved subjective outcomes for Task 4 (mean human score 3.53 vs. 3.13 for Task 3).

\subsection{Human Evaluation Analysis}
Human evaluation is synthetic but multi-rater and structured. For each sample, 10 participant scores are generated on a 1--5 scale from a proxy quality score plus participant bias and Gaussian noise. The proxy reward is:

\begin{equation}
r = 0.4\,S_{\text{pitch}} + 0.4\,D_{\text{rhythm}} + 0.2\,(1-R).
\end{equation}

This design keeps survey labels aligned with the optimization target while still producing inter-rater variance. At before-vs-after level, Task 4 improves mean human score (3.53 vs. 3.13) and slightly lowers average score spread (0.4120 vs. 0.4335), indicating better and more consistent outputs under the synthetic protocol.

Model-level survey means from the current export are: Task 1 (3.85), Markov baseline (3.75), Task 4 RLHF (3.53), Task 2 VAE samples (3.18), Task 2 interpolation (3.16), Task 3 Transformer (3.13) and Random baseline (3.10). These results show that RLHF improves the Task 3 policy in this setup, though absolute ranking remains sensitive to the reward design and dataset domain.

\section{Limitations and Future Work}
\subsection{Current Limitations}
This pipeline has several practical and methodological limitations. First, the event-token context cap (max sequence length 256) constrains long-form compositional planning; generated samples can still lose structure beyond local windows. Second, the RLHF stage is preference-aligned only with respect to a synthetic survey protocol, not fully real human listening sessions, so external validity remains limited. Third, evaluation remains symbolic-domain focused (piano-roll and token metrics), which does not directly capture timbral or expressive audio realism. Finally, model capacity is intentionally conservative for compute feasibility, which may underfit richer cross-genre dependencies.

\subsection{Future Directions}
The next logical step is to introduce real human-in-the-loop preference collection with active query selection, so annotation effort is focused on informative comparisons. On the modeling side, scaling Transformer capacity and context length, plus explicit hierarchical structure (phrase/section conditioning), should improve long-horizon coherence. Reward modeling can be strengthened with pairwise ranking objectives and uncertainty-aware calibration. Beyond symbolic MIDI generation, a hybrid symbolic-to-audio or end-to-end audio model (for example diffusion or neural codec language models) would broaden perceptual realism and practical music-production relevance.

\section{Conclusion}
This report documents a complete four-stage symbolic-music generation pipeline: LSTM autoencoder, VAE, Transformer and RLHF fine-tuning. The final report now integrates architecture/hyperparameter disclosure, expanded dataset-preprocessing detail, latent interpolation analysis and objective plus survey-based evaluation tables. The Transformer reached best validation perplexity 2.6312 after 50 cumulative epochs and the RLHF stage improved the Task 3 policy on rhythm diversity, repetition behavior and mean survey score (3.53 vs. 3.13). Together, these results provide a reproducible and empirically supported comparison across unsupervised and preference-aligned generation paradigms.

\nocite{*}

\bibliographystyle{IEEEtran}
\bibliography{references}

\end{document}
